% GENERATED FILE. Edit canonical lessons, then run npm run generate:books.
% canonical-source-hash: fnv1a64:2fd992417d5c376a
% canonical-lessons: PA-C271-uttar, PA-C271-dakkhan, PA-C271-purab, PA-C271-pachchham, PA-C271-ghumauna

\chapter{North, South, East, West}
\label{ch:pa-a2-north-south-east-west}
\hlchaptermodality{271}

\begin{chapteropening}
\textbf{By the end of this chapter:} \emph{I can say in which direction a place lies.}

Complete the last lesson of chapter 271: I can say in which direction a place lies.
\end{chapteropening}

\section[\texorpdfstring{\pa{ਉੱਤਰ} (uttar)}{ਉੱਤਰ (uttar)}]{\pa{ਉੱਤਰ} (uttar) — the north}

\label{lesson:PA-C271-uttar}



\begin{quote}
Before the new one: say the Punjabi for a mile, then the Punjabi for open ground.
\end{quote}

\subsection*{You'll want to know: \pa{ਉੱਤਰ}}
\textbf{\pa{ਉੱਤਰ}} — \emph{uttar} — \textquotedblleft{}the north\textquotedblright{}.

\begin{quote}
\pa{ਪਹਾੜ} \pa{ਪਿੰਡ} \pa{ਦੇ} \pa{ਉੱਤਰ} \pa{ਵਿੱਚ} \pa{ਹਨ।}
\end{quote}

\begin{quote}
\emph{pahāṛ piṇḍ de uttar vich han.}
\end{quote}

\begin{quote}
The mountains are to the north of the village.
\end{quote}

A masculine noun.

\begin{cousinweb}
\textbf{\pa{ਉੱਤਰ}} (\emph{uttar}) is inherited from Sanskrit \emph{uttara}, \textquotedblleft{}upper, northern\textquotedblright{}; CDIAL lists Punjabi \emph{uttar}, \textquotedblleft{}the north\textquotedblright{}, under it (CDIAL 1767).
\end{cousinweb}

\begin{grammarlens}[title={What you've built}]
One more word for this chapter.
\end{grammarlens}

\subsection*{Your turn}
\begin{itemize}
\raggedright
  \item \emph{Say it:} \emph{uttar}
  \item \emph{Say it:} \emph{uttar}, once more
  \item \emph{Say it:} \emph{pahāṛ piṇḍ de uttar vich han.} — The mountains are to the north of the village.
  \item \emph{Recall:} say the Punjabi for a mile, then the Punjabi for open ground, then say \emph{uttar} again
\end{itemize}

\begin{tcolorbox}[breakable,colback=teal!4,colframe=teal!35!black,title={Before you move on}]
What does \pa{ਉੱਤਰ} mean? (\textquotedblleft{}The north\textquotedblright{}.) Say it once more.
\end{tcolorbox}
\section[\texorpdfstring{\pa{ਦੱਖਣ} (dakkhaṇ)}{ਦੱਖਣ (dakkhaṇ)}]{\pa{ਦੱਖਣ} (dakkhaṇ) — the south}

\label{lesson:PA-C271-dakkhan}



\begin{quote}
Before the new one: say the Punjabi for open ground, then the Punjabi for north.
\end{quote}

\subsection*{You'll want to know: \pa{ਦੱਖਣ}}
\textbf{\pa{ਦੱਖਣ}} — \emph{dakkhaṇ} — \textquotedblleft{}the south\textquotedblright{}.

\begin{quote}
\pa{ਸਮੁੰਦਰ} \pa{ਦੱਖਣ} \pa{ਵਿੱਚ} \pa{ਹੈ।}
\end{quote}

\begin{quote}
\emph{samundar dakkhaṇ vich hai.}
\end{quote}

\begin{quote}
The sea is in the south.
\end{quote}

A masculine noun.

\begin{cousinweb}
\textbf{\pa{ਦੱਖਣ}} (\emph{dakkhaṇ}) is inherited from Sanskrit \emph{dakṣiṇa}, \textquotedblleft{}right (hand), southern\textquotedblright{}; CDIAL lists Punjabi \emph{dakkhaṇ}, \textquotedblleft{}the south\textquotedblright{}, under it (CDIAL 6119).
\end{cousinweb}

\begin{grammarlens}[title={What you've built}]
One more word for this chapter.
\end{grammarlens}

\subsection*{Your turn}
\begin{itemize}
\raggedright
  \item \emph{Say it:} \emph{dakkhaṇ}
  \item \emph{Say it:} \emph{dakkhaṇ}, once more
  \item \emph{Say it:} \emph{samundar dakkhaṇ vich hai.} — The sea is in the south.
  \item \emph{Recall:} say the Punjabi for open ground, then the Punjabi for north, then say \emph{dakkhaṇ} again
  \item \emph{Write it:} copy \textbf{\pa{ਦੱਖਣ}} once, checking the whole word
\end{itemize}

\begin{tcolorbox}[breakable,colback=teal!4,colframe=teal!35!black,title={Before you move on}]
What does \pa{ਦੱਖਣ} mean? (\textquotedblleft{}The south\textquotedblright{}.) Say it once more.
\end{tcolorbox}
\section[\texorpdfstring{\pa{ਪੂਰਬ} (pūrab)}{ਪੂਰਬ (pūrab)}]{\pa{ਪੂਰਬ} (pūrab) — the east}

\label{lesson:PA-C271-purab}



\begin{quote}
Before the new one: say the Punjabi for north, then the Punjabi for south.
\end{quote}

\subsection*{You'll want to know: \pa{ਪੂਰਬ}}
\textbf{\pa{ਪੂਰਬ}} — \emph{pūrab} — \textquotedblleft{}the east\textquotedblright{}.

\begin{quote}
\pa{ਸੂਰਜ} \pa{ਪੂਰਬ} \pa{ਵਿੱਚ} \pa{ਚੜ੍ਹਦਾ} \pa{ਹੈ।}
\end{quote}

\begin{quote}
\emph{sūraj pūrab vich chaṛhdā hai.}
\end{quote}

\begin{quote}
The sun rises in the east.
\end{quote}

A masculine noun.

\begin{cousinweb}
\textbf{\pa{ਪੂਰਬ}} (\emph{pūrab}) goes back to Sanskrit \emph{pūrva}, \textquotedblleft{}in front; eastern\textquotedblright{} (Monier-Williams), in a half-learned form reshaped to Punjabi sounds; CDIAL gives no Punjabi form under it (CDIAL 8343).
\end{cousinweb}

\begin{grammarlens}[title={What you've built}]
One more word for this chapter.
\end{grammarlens}

\subsection*{Your turn}
\begin{itemize}
\raggedright
  \item \emph{Say it:} \emph{pūrab}
  \item \emph{Say it:} \emph{pūrab}, once more
  \item \emph{Say it:} \emph{sūraj pūrab vich chaṛhdā hai.} — The sun rises in the east.
  \item \emph{Recall:} say the Punjabi for north, then the Punjabi for south, then say \emph{pūrab} again
  \item \emph{Write it:} copy \textbf{\pa{ਪੂਰਬ}} once, checking the whole word
\end{itemize}

\begin{tcolorbox}[breakable,colback=teal!4,colframe=teal!35!black,title={Before you move on}]
What does \pa{ਪੂਰਬ} mean? (\textquotedblleft{}The east\textquotedblright{}.) Say it once more.
\end{tcolorbox}
\section[\texorpdfstring{\pa{ਪੱਛਮ} (pachchham)}{ਪੱਛਮ (pachchham)}]{\pa{ਪੱਛਮ} (pachchham) — the west}

\label{lesson:PA-C271-pachchham}



\begin{quote}
Before the new one: say the Punjabi for south, then the Punjabi for east.
\end{quote}

\subsection*{You'll want to know: \pa{ਪੱਛਮ}}
\textbf{\pa{ਪੱਛਮ}} — \emph{pachchham} — \textquotedblleft{}the west\textquotedblright{}.

\begin{quote}
\pa{ਬੱਸ} \pa{ਅੱਡਾ} \pa{ਸ਼ਹਿਰ} \pa{ਦੇ} \pa{ਪੱਛਮ} \pa{ਵਿੱਚ} \pa{ਹੈ।}
\end{quote}

\begin{quote}
\emph{bass aḍḍā shahir de pachchham vich hai.}
\end{quote}

\begin{quote}
The bus stand is in the west of the town.
\end{quote}

A masculine noun.

\begin{cousinweb}
\textbf{\pa{ਪੱਛਮ}} (\emph{pachchham}) goes back to Sanskrit \emph{paścima}, \textquotedblleft{}behind; western\textquotedblright{} (Monier-Williams), in a half-learned form reshaped to Punjabi sounds; CDIAL lists only an older Punjabi \emph{pachvā}, \textquotedblleft{}western\textquotedblright{}, under it (CDIAL 8007).
\end{cousinweb}

\begin{grammarlens}[title={What you've built}]
One more word for this chapter.
\end{grammarlens}

\subsection*{Your turn}
\begin{itemize}
\raggedright
  \item \emph{Say it:} \emph{pachchham}
  \item \emph{Say it:} \emph{pachchham}, once more
  \item \emph{Say it:} \emph{bass aḍḍā shahir de pachchham vich hai.} — The bus stand is in the west of the town.
  \item \emph{Recall:} say the Punjabi for south, then the Punjabi for east, then say \emph{pachchham} again
  \item \emph{Write it:} copy \textbf{\pa{ਪੱਛਮ}} once, checking the whole word
\end{itemize}

\begin{tcolorbox}[breakable,colback=teal!4,colframe=teal!35!black,title={Before you move on}]
What does \pa{ਪੱਛਮ} mean? (\textquotedblleft{}The west\textquotedblright{}.) Say it once more.
\end{tcolorbox}
\section[\texorpdfstring{\pa{ਘੁਮਾਉਣਾ} (ghumāuṇā)}{ਘੁਮਾਉਣਾ (ghumāuṇā)}]{\pa{ਘੁਮਾਉਣਾ} (ghumāuṇā) — to show someone around; to turn round}

\label{lesson:PA-C271-ghumauna}



\begin{quote}
Before the new one: say the Punjabi for east, then the Punjabi for west.
\end{quote}

\subsection*{You'll want to know: \pa{ਘੁਮਾਉਣਾ}}
\textbf{\pa{ਘੁਮਾਉਣਾ}} — \emph{ghumāuṇā} — \textquotedblleft{}to show someone around; to turn round\textquotedblright{}.

\begin{quote}
\pa{ਮੈਂ} \pa{ਮਹਿਮਾਨਾਂ} \pa{ਨੂੰ} \pa{ਸ਼ਹਿਰ} \pa{ਘੁਮਾਵਾਂਗਾ।}
\end{quote}

\begin{quote}
\emph{maiṁ mahimānāṁ nū̃ shahir ghumāvāṅgā.}
\end{quote}

\begin{quote}
I will show the guests around the town.
\end{quote}

Past (I …): \pa{ਮੈਂ} \pa{ਮਨਨ} \pa{ਨੂੰ} \pa{ਸ਼ਹਿਰ} \pa{ਘੁਮਾਇਆ}, \emph{maiṁ Manan nū̃ shahir ghumāiā}.

\begin{cousinweb}
\textbf{\pa{ਘੁਮਾਉਣਾ}} (\emph{ghumāuṇā}) is the causative partner of \emph{ghummaṇā}, \textquotedblleft{}to turn, wander\textquotedblright{}; CDIAL lists Punjabi \emph{ghummaṇā} under a reconstructed \emph{*ghummati}, \textquotedblleft{}revolves\textquotedblright{} (CDIAL 4485).
\end{cousinweb}

\begin{grammarlens}[title={What you've built}]
One more word for this chapter.
\end{grammarlens}

\subsection*{Your turn}
\begin{itemize}
\raggedright
  \item \emph{Say it:} \emph{ghumāuṇā}
  \item \emph{Say it:} \emph{ghumāuṇā}, once more
  \item \emph{Say it:} \emph{maiṁ mahimānāṁ nū̃ shahir ghumāvāṅgā.} — I will show the guests around the town.
  \item \emph{Recall:} say the Punjabi for east, then the Punjabi for west, then say \emph{ghumāuṇā} again
  \item \emph{Write it:} copy \textbf{\pa{ਘੁਮਾਉਣਾ}} once, checking the whole word
\end{itemize}

\begin{tcolorbox}[breakable,colback=teal!4,colframe=teal!35!black,title={Before you move on}]
What does \pa{ਘੁਮਾਉਣਾ} mean? (\textquotedblleft{}To show someone around; to turn round\textquotedblright{}.) Say it once more.
\end{tcolorbox}
